\section{\texorpdfstring{$\Ham(S^2)$}{Ham(S2)} 上的\personPolterovich{}定理}
\label{sec:polterovich-Ham-S2}

联系\personHofer{}距离，一个自然问题是：对给定紧辛流形，$\Ham(M,\omega)$ 的直径究竟
有限还是无限？

\begin{definition}[\personHofer{}直径]
\label{def:hofer-diameter}
设 $(M,\omega)$ 是辛流形，$\Ham(M,\omega)$ 是其哈密顿微分同胚群。（若
$(M,\omega)$ 是开的，则改取 $\Ham_c(M,\omega)$。）定义
\[
 \operatorname{diam}(\Ham(M,\omega))
 =\sup_{\phi\in\Ham(M,\omega)}\|\phi\|.
\]
\end{definition}

\stfhIndex{I154}{Boerjieluoweiqi}{波尔捷罗维奇（Polterovich）}
\personOstrover{} (Os03) 与\personPolterovich{} (Po98b) 分别证明：在
$\pi_2(M)=\{0\}$ 的非球面情形以及 $S^2$ 的情形，\personHofer{}直径为无穷。

\begin{theorem}[\personOstrover{}]
\label{thm:ostrover-hofer-infinite}
设 $(M,\omega)$ 是满足 $\pi_2(M)=0$ 的闭辛流形。则
\stfhIndex{I111}{Huofeierzhijing}{霍费尔直径（Hofer's diameter）}
$\Ham(M,\omega)$ 的\termHoferDiameter{}为无穷。
\end{theorem}

\personUsher{}在 (Ush10a) 中给出了这一结果的各种推广。另一方面，\personPolterovich{}证明了：

\begin{theorem}[\personPolterovich{}]
\label{thm:polterovich-hofer-S2-infinite}
$\Ham(S^2)$ 的\termHoferDiameter{}为无穷。
\end{theorem}

\personOstrover{}的定理使用哈密顿\personFloer{}同调证明，而\personPolterovich{}使用拉格朗日交\personFloer{}同调。我们将在第四部分回到\personOstrover{}的证明。本节紧随\personPolterovich{}的原始论证，
给出他的证明。

先从应用于哈密顿环路的拉格朗日悬垂构造开始。考虑满足
\[
 H:[0,1]\times M\to\mathbb R,
 \qquad \phi_H^1=\mathrm{id}
\]
的哈密顿量。考虑任意闭辛流形 $(M,\omega)$。设
$\phi\in\Ham(M,\omega)$，$F$ 是满足 $F\mapsto\phi$ 的归一化哈密顿量。

\begin{definition}
\label{def:periodic-hamiltonian-loops}
以 $\mathcal H_{\mathrm{loop}}$ 表示生成哈密顿环路的一周期哈密顿量 $H$（即
$H(t,x)=H(t+1,x)$）的集合；换言之，$\phi_H$ 满足
$\phi_H^1=\mathrm{id}$。
\end{definition}

回忆，对任意给定的哈密顿环路，总能重新参数化它，使
$\phi_H^t\equiv\mathrm{id}$ 在 $t=0,1$ 附近成立，从而使相伴哈密顿量 $H$ 为
一周期的。特别地，对某个常数 $c$，在
$[0,\delta]\cup[1-\delta,1]\times M$ 上 $H\equiv c$。

对定义 3.7.6 中哈密顿同痕 $\phi_H$，考虑赤道
$L\subset(M,\omega)$ 的悬垂，记作
\[
 N(L,\phi_H)\subset T^*L\times M,
\]
并令 $N_0=N(L,\mathrm{id})$。注意 $N_0$ 就是积 $o_L\times L$，且命题 3.7.7
蕴含 $\mathcal L=N(L,\phi_H)$ 与 $N_0$ 哈密顿同痕。此外，因为
$L\subset(M,\omega)$ 单调且最小马斯洛夫数为 2，$N_0$ 亦然。所以
$N_0$ 的\personFloer{}同调 $HF(N_0,N(L,\phi_H))$ 在\personNovikov{}环
$\Lambda^{\mathbb Z_2}$ 上良定义，并由定理 16.4.10 同构于 $HF(N_0,N_0)$。
另一方面，我们有如下计算。

\begin{lemma}
\label{lem:floer-product-suspension}
假设 $L\subset(M,\omega)$ 单调且 $\Sigma_L\geq2$。则
\[
 HF(N_0,N_0)\cong HF(o_L,o_L)\otimes HF(L,L;M).
\]
\end{lemma}

\begin{proof}
由于 $N_0$ 单调且最小马斯洛夫数至少为 2，$N_0$ 与其哈密顿形变的相伴\personFloer{}同调
不依赖于同痕和殆复结构的选择（再用定理 16.4.10）。
\footnote{译者注：本引理只假设 $\Sigma_L\geq2$，而原书断言积拉格朗日
$N_0=o_L\times L$ 的最小马斯洛夫数恒为 2；一般只能由积结构得到至少为 2。
译文按证明实际所需的下界改正。}

注意拉格朗日子流形 $N_0=o_L\times L$ 是积辛流形 $T^*L\times M$ 中的直积。
利用 $HF$ 在哈密顿同痕下的不变性（仍用定理 16.4.10），可在
$T^*L\times M$ 上考虑积殆复结构 $J_0$，以及 $N_0$ 的通有积哈密顿同痕
$\phi_{H'}=\phi_{H_1'}\times\phi_{H_2'}$。相伴\personFloer{}流线方程分裂，因此相应\personFloer{}复形 $(CF(N_0,\phi_{H'}^1(N_0)),\partial)$ 成为
\[
 CF(N_0,\phi_{H'}^1(N_0))=CF(o_L,\phi_{H_1'}^1(o_L))
 \otimes CF(L,\phi_{H_2'}^1(L)),
\]
\footnote{译者注：原书在本式左端突然写未定义的 $L'$；由上一句正在讨论的\personFloer{}复形可唯一确定其应为 $\phi_{H'}^1(N_0)$，译文据此补正。}
其边界映射为
\[
 \partial(x\otimes y)=\partial(x)\otimes y+x\otimes\partial(y).
\]
于是立即得到同构
\[
 HF(N_0,N_0)\cong HF(o_L,o_L)\otimes HF(L,L;M).
\]
\footnote{译者注：引理陈述给出张量积 $\otimes$，且证明中的链复形也是张量积，
但原书证明末句却印作直和 $\oplus$；译文据引理陈述及前一公式改正。}
\end{proof}

\stfhIndex{I181}{Wendingxiangjiaoxingzhi}{稳定相交性质（stable intersection property）}
直接推论是：若 $L\subset M$ 满足 $HF(L,L;M)\neq0$，则它具有如下稳定相交性质。

\begin{corollary}
\label{cor:stable-intersection-from-HF}
假设 $L\subset(M,\omega)$ 单调且 $\Sigma_L\geq2$，因而\personFloer{}同调
$HF(L,L;M)$ 有定义。再假设 $HF(L,L;M)\neq0$。则对 $(M,\omega)$ 的任意哈密顿
环路 $\phi_H$，都有
\[
 (o_L\times L)\cap N(L,\phi_H)\neq\varnothing.
\]
\end{corollary}

从下面的定义看，这样的拉格朗日子流形 $L$ 具有稳定拉格朗日交性质。

\begin{definition}
\label{def:stable-lagrangian-intersection}
若对任意哈密顿环路 $\phi_H$（即满足 $\phi_H^1=\mathrm{id}$），都有
$(o_L\times L)\cap N(L,\phi_H)\neq\varnothing$，则称拉格朗日子流形
$L\subset(M,\omega)$ 具有\termStableLagrangianIntersectionProperty{}。
\end{definition}

如下\personFloer{}同调非消失结果实质上已在 (Oh93b) 中证明。

\begin{theorem}
\label{thm:equator-floer-nonzero}
考虑 $S_{\mathrm{eq}}^1\subset S^2$。则
\[
 HF(S_{\mathrm{eq}}^1,S_{\mathrm{eq}}^1;S^2)
 \cong H(S^1,\Lambda^{\mathbb Z_2}).
\]
特别地，$S_{\mathrm{eq}}^1$ 具有稳定拉格朗日交性质。
\end{theorem}

\begin{proof}
在 $S_{\mathrm{eq}}^1$ 上取一个 $C^2$-小\personMorse{}函数 $f$；它恰在赤道的一对对跖点
有两个临界点，并生成给定赤道的一个小旋转。
\footnote{译者注：原书把这两个位于赤道上的临界点称为 $S^2$ 的南极和北极，但
南北极不在 $S_{\mathrm{eq}}^1$ 上；证明只需赤道上的一对对跖临界点，译文据此改正。}
应用\personRiemann{}映射定理可直接检验，在本情形
中厚--薄分解给出
\[
 \partial=\partial_f^{\mathrm{Morse}}+\partial',
\]
其中 $\partial'$ 由两条\personFloer{}流线贡献：从上下半球切掉
$S_{\mathrm{eq}}^1$ 与 $\phi_f^1(S_{\mathrm{eq}}^1)$ 之间的“薄”流线，就得到它们。
恰有两条这样的“厚”流线。因此矩阵系数为
\[
 T^{\omega(u_1)}+T^{\omega(u_2)},
\]
其中 $\omega(u_i)$ 是相应\personFloer{}流线 $u_i$ 的面积。另一方面，由
$S_{\mathrm{eq}}^1$ 的单调性，两个面积相同，故矩阵系数变为 $2T^\lambda$，其中
$\lambda$ 是公共面积。因此在 $\Lambda^{\mathbb Z_2}$ 中该矩阵系数为零。这说明
$\partial=\partial_f^{\mathrm{Morse}}$，证明完毕。
\end{proof}

接下来回忆拉格朗日悬垂的显式定义
\[
 \iota_{\phi_H,L}:(t,x)\longmapsto
 (t,-H(t,\phi_H^t(x)),\phi_H^t(x)),
\]
其像就是 $N(L,\phi_H)$。交点
\[
 (t,a,y)\in N_0\cap N(L,\phi_H)
\]
由方程
\begin{equation}
 0=-H(t,y)=a,
 \qquad y\in\operatorname{Fix}\phi_H^t\cap L
\label{eq:17-suspension-intersection-characterization}
\end{equation}
刻画。

还回忆，在 $[-\delta,\delta]\times M$ 上 $H\equiv c$，其中把
$[-\delta,\delta]$ 看作 $S^1=\mathbb R/\mathbb Z$ 的子集。为去掉平凡交点，选取
可任意小的 $c\neq0$，再重新归一化 $H_t$，使其满足归一化条件。做到这一点时可让
$\|H\|$ 任意小。

有了这些准备，可以完成定理 17.6.3 的证明。

\begin{proof}[定理 17.6.3 的证明]
先用周期哈密顿量集合 $\mathcal H_\phi$（见定义 17.6.4）描述
$\phi\in\Ham(M,\omega)$ 的\personHofer{}范数 $\|\phi\|$。

\begin{lemma}
\label{lem:hofer-norm-loop-infimum}
设 $\phi\in\Ham(M,\omega)$ 且 $\phi=\phi_F^1$。则
\[
 \|\phi\|=\inf_{H\in\mathcal H_{\mathrm{loop}}}\|F-H\|.
\]
\end{lemma}

\begin{proof}
证明留作习题。
\end{proof}

考虑哈密顿同痕 $\phi_H$ 所给赤道
$L=S_{\mathrm{eq}}^1\subset S^2$ 的悬垂，记为
\[
 N(S_{\mathrm{eq}}^1,\phi_H)\subset T^*S^1\times S^2,
\]
其中
\[
 S^1=S^1(2)=\mathbb R/2\mathbb Z,
\]
并令 $N_0=N(S_{\mathrm{eq}}^1,\mathrm{id})$。
\footnote{译者注：原书此式在商号后印有额外的 $c$，写作
$\mathbb R/c\,2\mathbb Z$；此前的 $c$ 是可任意小的哈密顿常数，不可能定义固定的
$S^1(2)$，故译文改为 $\mathbb R/2\mathbb Z$。}
注意 $N_0$ 就是积 $o_{T^*S_{\mathrm{eq}}^1}\times S_{\mathrm{eq}}^1$；命题 3.7.7
蕴含，只要哈密顿环路在 $\Ham(S^2,\omega)$ 中可缩，
$\mathcal L=N(S_{\mathrm{eq}}^1,\phi_H)$ 就与 $N_0$ 哈密顿同痕。

另一方面，$\pi_1(\Ham(S^2))\cong\mathbb Z_2$，且沿垂直于给定赤道
$S_{\mathrm{eq}}^1$ 的轴作整周旋转生成一个非平凡环路。这个环路把
$S_{\mathrm{eq}}^1$ 映到自身，且相伴哈密顿量在 $S_{\mathrm{eq}}^1$ 上消失。因此
相伴悬垂仍为 $N_0$。所以，对所有哈密顿环路 $\phi_H$，
$N(S_{\mathrm{eq}}^1,\phi_H)$ 都与 $N_0$ 哈密顿同痕。

又因 $S_{\mathrm{eq}}^1\subset S^2$ 单调且最小马斯洛夫数为 2，$N_0$ 亦然。因此
$N_0$ 的\personFloer{}同调
\[
 HF(N_0,N(S^1,\phi_H))
\]
在\personNovikov{}环 $\Lambda^{\mathbb Z_2}$ 上良定义，并由定理 16.4.11 同构于
$HF(N_0,N_0)$。
然后对
\[
 HF(o_{T^*S^1},o_{T^*S^1})\cong H(S^1)
 \cong\Lambda^{\mathbb Z_2}\oplus\Lambda^{\mathbb Z_2}
\]
应用余切丛同构，并对
\[
 HF(S_{\mathrm{eq}}^1,S_{\mathrm{eq}}^1;S^2)
 \cong H(S_{\mathrm{eq}}^1)
 \cong\Lambda^{\mathbb Z_2}\oplus\Lambda^{\mathbb Z_2}\neq0
\]
应用定理 17.6.8。由推论 17.6.6，$S_{\mathrm{eq}}^1$ 具有稳定拉格朗日交性质。

于是由~\eqref{eq:17-suspension-intersection-characterization}，存在
$(t,y)\in S^1\times L$，使得 $H(t,y)=0$。注意上述证明依赖于
$\Ham(S^2)$ 的基本群有限这一事实。它与引理 17.6.9 立即给出如下推论。

\begin{corollary}
\label{cor:hofer-lower-bound-equator}
设归一化哈密顿量 $F$ 满足
$F|_{S_{\mathrm{eq}}^1}\equiv c>0$，并令 $\phi=\phi_F^1$。则
$\|\phi\|\geq c$。
\end{corollary}

最后，在此推论中可选择归一化哈密顿量 $F$，使 $c$ 任意大，证明即告完成。
\end{proof}

事实上，只要任意闭 $(M,\omega)$ 包含一个具有稳定相交性质的紧拉格朗日子流形，
上述证明就适用。

\begin{theorem}
\label{thm:stable-intersection-implies-infinite-\personHofer{}}
假设 $(M,\omega)$ 包含一个紧拉格朗日子流形 $L$，并且 $L$ 具有\termStableLagrangianIntersectionProperty{}。则
$\operatorname{diam}\Ham(M,\omega)$ 为无穷。
\end{theorem}

在这方面，我们已经在上述证明中说明
$S_{\mathrm{eq}}^1\subset S^2$ 具有稳定拉格朗日交性质。用同一论证还可推出，
$\mathbb{RP}^n$ 以及 $\mathbb{CP}^n$ 中的\personClifford{}环面 $T^n$ 也具有这一性质。

\begin{remark}
\label{rem:hofer-versus-spectral-diameter}
定理 17.6.3 与将在第 22.3 节介绍的谱直径有限性定理形成鲜明对比。我们将在第
22.5 节证明
\[
 \operatorname{diam}_\rho(\Ham(S^2))\leq4\pi
 \quad(=S^2\text{ 的面积}).
\]
\end{remark}
